\begin{tikzpicture}[scale=0.88,>=Latex,line cap=round,line join=round]
  \begin{scope}
    \node[font=\small] at (2.2,2.75) {限流接法};

    % Source and return wire. The source branch is deliberately broken
    % between the two plates so it cannot be mistaken for a short circuit.
    \draw[thick] (0,0) -- (4.4,0) -- (4.4,2.0) -- (3.75,2.0);
    \draw[thick] (0,0) -- (0,0.62);
    \draw[thick] (-0.17,0.62) -- (0.17,0.62);
    \draw[thick] (-0.28,0.92) -- (0.28,0.92);
    \draw[thick] (0,0.92) -- (0,2.62) -- (1.38,2.62);
    \node[font=\scriptsize] at (-0.42,0.77) {$E$};

    % Rheostat: the slider and one end are the two terminals in use.
    \draw[thick] (0.78,1.75) rectangle (2.18,2.25);
    \node at (1.48,2.0) {$R$};
    \draw[thick,->,blue] (1.38,2.62) -- (1.48,2.25);
    \node[font=\scriptsize,blue] at (0.92,2.42) {滑片};
    \draw[thick] (2.18,2.0) -- (2.62,2.0);

    % Load.
    \draw[thick] (2.62,1.72) rectangle (3.42,2.28);
    \node at (3.02,2.0) {$R_x$};
    \draw[thick] (3.42,2.0) -- (3.75,2.0);
    \draw[->,red,thick] (3.75,2.36) -- (3.18,2.36);
    \node[font=\scriptsize,red] at (2.62,2.42) {$I$ 可调};

    \node[font=\scriptsize,align=center] at (2.2,-0.55)
      {滑片与电阻丝一端接入，和 $R_x$ 串联\\电流通常不能从零连续调起};
  \end{scope}

  \begin{scope}[xshift=6.3cm]
    \node[font=\small] at (2.2,2.75) {分压接法};

    % Source. Keep an actual gap between the cell plates.
    \draw[thick] (0,0) -- (0,0.62);
    \draw[thick] (-0.17,0.62) -- (0.17,0.62);
    \draw[thick] (-0.28,0.92) -- (0.28,0.92);
    \draw[thick] (0,0.92) -- (0,2.0) -- (0.78,2.0);
    \draw[thick] (0,0) -- (4.4,0) -- (4.4,2.0) -- (3.62,2.0);
    \node[font=\scriptsize] at (-0.42,0.77) {$E$};

    % Entire resistance wire is connected across the source.
    \draw[thick] (0.78,1.75) rectangle (3.62,2.25);
    \node at (2.20,2.0) {$R$};

    % Slider feeds the load; the other load terminal returns to one end.
    \coordinate (S) at (2.05,2.25);
    \draw[thick,->,blue] (2.05,2.72) -- (S);
    \node[font=\scriptsize,blue] at (1.55,2.46) {滑片};
    \draw[thick] (S) -- (2.05,1.15) -- (2.78,1.15);
    \draw[thick] (2.78,0.82) rectangle (3.58,1.48);
    \node at (3.18,1.15) {$R_x$};
    \draw[thick] (3.58,1.15) -- (4.4,1.15);
    \draw[->,red,thick] (2.25,1.48) -- (2.72,1.48);
    \node[font=\scriptsize,red] at (2.18,1.02) {$U_x$};

    \node[font=\scriptsize,align=center] at (2.2,-0.55)
      {电阻丝两端跨接电源，滑片给 $R_x$ 分压\\电压可从近似零连续调起};
  \end{scope}
\end{tikzpicture}
